\begin{abstract}
We present TransitDuet, a bi-level reinforcement learning framework that jointly optimizes bus timetable planning and real-time holding control.
This joint problem is fundamentally harder than either task in isolation: the planning and control layers operate on different timescales, observe different state spaces, and interact through delayed, asynchronous feedback---properties that violate the synchronous co-execution assumption underlying existing hierarchical reinforcement learning methods.
TransitDuet addresses this through three mechanisms.
First, a goal-conditioned cross-level coupling in the spirit of HIRO: the upper outputs a per-dispatch target-headway shift, and the lower's Lagrangian cost is computed against the resulting goal headway, with no upper-advantage injection into the lower's transitions. This avoids the covariate shift that destabilises a literal cross-advantage augmentation while keeping the upper--lower interaction unidirectional and stable.
Second, a measurement-projection scheme based on online gradient descent enforces fleet-size constraints without reward shaping.
Third, Lagrangian relaxation within the lower policy penalizes headway deviations, converting a hard service-regularity requirement into a soft dual objective.
We evaluate TransitDuet in a calibrated microsimulation of a high-frequency bus corridor with stochastic demand and traffic.
A complementary technique, target-policy-corrected lower SAC (TPC), uses importance-reweighting to align the lower's training distribution with its eventual deployment distribution; combined with the goal-conditioned coupling and trained for 300 episodes per seed, TransitDuet attains the best mean on every metric we report. Across three random seeds and an elastic fleet budget sampled per episode from $[8, 16]$, TransitDuet reduces composite cost by $6.5\%$ ($1.444$ vs $1.544$) and overshoot by $10.4\%$ relative to the strongest hand-tuned fixed timetable, while still improving wait, CV, and overshoot over GA and CMA-ES under the same 300-episode training and held-out evaluation protocol. Seed variance on composite cost is $1.7$--$2.5\times$ smaller than the search baselines, and the learned policy produces a Pareto frontier across nine fleet sizes from a single training run. Generalisation degrades gracefully: wait rises only $25\%$ when route variance is doubled beyond the training distribution.
\end{abstract}
